% Evans 7.5: printed pp.420--423 / English PDF pp.437--440.
% Chinese: printed pp.297--299 / scan pp.315--317.
\Needspace{60mm}
\Needspace{55mm}
\section{Problems and Worked Solutions}
\label{sec:7-5}
The following are concise restatements of all 16 source exercises, with
editorial solutions. Unless specified otherwise, $U$ is bounded with smooth
boundary and all classical solutions have enough regularity for the indicated
integrations. The source problem numbers are retained in the paragraph labels.

\Needspace{30mm}
\paragraph{1. Neumann heat uniqueness.}
Prove uniqueness for $u_t-\Delta u=f$, $\partial_\nu u=0$, $u(0)=g$.
\begin{proof}[Solution]
A difference $w$ has zero initial value and homogeneous equation and boundary
data. Multiplication by $w$ yields
$\frac12\frac d{dt}\|w\|_2^2+\|Dw\|_2^2=0$; the boundary term is
$\int_{\partial U}w\partial_\nu w=0$. Thus $w=0$. Neumann constants do
not obstruct evolutionary uniqueness once the initial value is fixed.
\end{proof}

\Needspace{30mm}
\paragraph{2. Sharp Dirichlet heat decay.}
Show $\|u(t)\|_2\leq e^{-\lambda_1t}\|g\|_2$ for homogeneous heat flow.
\begin{proof}[Solution]
The energy identity and the Rayleigh characterization from
\S6.5 give
$\frac d{dt}\|u\|_2^2=-2\|Du\|_2^2\leq-2\lambda_1\|u\|_2^2$.
Multiply by $e^{2\lambda_1t}$, integrate, then take square roots.
Equality holds for a first Dirichlet eigenfunction; this fixes the factor two.
\end{proof}
\Needspace{55mm}
\PDEAnnotation{平方范数的衰减率}{$\|u\|_2^2$ 以 $2\lambda_1$ 衰减，而 $\|u\|_2$ 以 $\lambda_1$ 衰减。Neumann 问题有常数模态，不能直接使用正的 Dirichlet 谱隙。}

\Needspace{30mm}
\paragraph{3. Adjoint dynamics.}
For $u_t+Lu=0$, $u(0)=g$, and $v_t-L^*v=0$, $v(T)=h$, both with
zero Dirichlet data, prove $(v(0),g)=(u(T),h)$.
\begin{proof}[Solution]
The formal adjoint is defined by $(Lu,v)=(u,L^*v)$ on the zero-boundary
smooth functions, including the coefficient derivative in the adjoint drift.
Therefore $\frac d{dt}(u,v)=-(Lu,v)+(u,L^*v)=0$.
Evaluate at $0,T$. The backward adjoint problem is prescribed terminal data;
it is not a forward heat problem with the diffusion sign reversed arbitrarily.
\end{proof}

\Needspace{30mm}
\paragraph{4. Galerkin approximation for Poisson's equation.}
For nested spans of a complete $V$-orthogonal basis, let
$\int Du_m\cdot Dw_k=\int fw_k$ for $k\leq m$. Show weak convergence
to the zero-Dirichlet Poisson solution.
\begin{proof}[Solution]
Testing with $u_m$ and Poincar\'e bound $\|u_m\|_V$ uniformly.
Each weak subsequential limit satisfies the equation on every fixed span,
hence on $V$ by density. Uniqueness identifies every subsequence.
In fact one has stronger convergence: if $u$ is the weak solution,
$\int D(u-u_m)\cdot Dv_m=0$, so
\[
 \|D(u-u_m)\|_2\leq\inf_{v_m\in V_m}\|D(u-v_m)\|_2\longrightarrow0.
\]
The finite systems are invertible since their gradient Gram matrices are
positive definite. No $L^2$ orthonormality of this exercise's basis is needed.
\end{proof}

\Needspace{30mm}
\paragraph{5. Identifying a weak time derivative.}
If $u_k\rightharpoonup u$ in $L^2V$ and $u_k'\rightharpoonup z$ in
$L^2V'$, show $z=u'$.
\begin{proof}[Solution]
For $\phi\in C_c^1(0,T)$, $w\in V$,
$\int\langle u_k',w\rangle\phi=-\int(u_k,w)\phi'$.
Both sides converge by the stipulated weak convergences. The resulting
identity is exactly the definition of the $V'$-valued distributional derivative.
\end{proof}

\Needspace{30mm}
\paragraph{6. The essential supremum survives a weak limit.}
If $u_k\rightharpoonup u$ in $L^2(0,T;H)$ and
$\|u_k(t)\|_H\leq C$ a.e., show $\|u(t)\|_H\leq C$ a.e.
\begin{proof}[Solution]
For any measurable $E\subset(0,T)$, $\mathbf1_Eu\in L^2(0,T;H)$ is a
valid fixed weak test. Thus
\[
 \int_E\|u\|_H^2=\lim_k\int_E(u_k,u)\leq C\int_E\|u\|_H.
\]
Choosing $E=\{\|u\|_H>C+\delta\}$ makes its integrand difference at least
$\delta\|u\|_H$, so $E$ has measure zero. Take $\delta=1/j$.
This argument works even without assuming $H$ separable.
\end{proof}
\Needspace{55mm}
\PDEAnnotation{一致上界与弱极限}{逐时刻弱收敛并非已知。用固定测试函数 $\mathbf1_Eu$ 把问题化为每个时间集合上的积分不等式，直接得到极限的本质上确界。}

\Needspace{30mm}
\paragraph{7. Pointwise decay from a positive potential.}
For $u_t-\Delta u+cu=0$, zero Dirichlet data and $c\geq\gamma>0$, prove
$|u(x,t)|\leq C e^{-\gamma t}$.
\begin{proof}[Solution]
Set $z=e^{\gamma t}u$. Then $z_t-\Delta z+(c-\gamma)z=0$ and the
potential is nonnegative. Comparison with the constants
$\pm\|g\|_\infty$ yields $|z|\leq\|g\|_\infty$ on every finite cylinder.
The constant is independent of its final time, proving the assertion.
\end{proof}

\Needspace{30mm}
\paragraph{8. Positivity with a bounded signed potential.}
For the same equation, $g\geq0$ and bounded $c$, show $u\geq0$.
\begin{proof}[Solution]
Choose $\lambda\geq\max(0,-\inf c)$ and $z=e^{-\lambda t}u$.
Then $z_t-\Delta z+(c+\lambda)z=0$. The minimum principle gives $z\geq0$,
and multiplication by $e^{\lambda t}>0$ proves the claim.
\end{proof}

\Needspace{30mm}
\paragraph{9. The second-order coercive estimate, including its boundary terms.}
Prove the source inequalities (54) in \S7.1.3 and (59) in \S7.2.3:
\begin{equation}\label{eq:7-h2-coercive}
 \beta\|u\|_{H^2}^2\leq(Lu,-\Delta u)+\gamma\|u\|_2^2
 \qquad(u\in H^2(U)\cap V).
\end{equation}
\begin{proof}[Solution]
Take first smooth $u=0$ on $\partial U$ and smooth symmetric $A$.
Expand the divergence operator as $Lu=-a^{ij}u_{ij}+\widetilde b^iu_i+cu$.
For the principal term put $I=\int_Ua^{ij}u_{ij}u_{kk}$. Integrate first
in $k$ and then in $j$:
\[
 \begin{split}
 I={}&\int_U a^{ij}u_{ik}u_{jk}
 +\int_U\big((\partial_j a^{ij})u_{ik}u_k
                 -(\partial_k a^{ij})u_{ij}u_k\big)\\
 &+\int_{\partial U}\big((A:D^2u)\partial_\nu u
                         -(D^2u\,A\nu)\cdot Du\big).
 \end{split}
\]
The volume errors are bounded by
$\delta\|D^2u\|_2^2+C_\delta\|Du\|_2^2$.
The first volume term is at least $\theta\|D^2u\|_2^2$.
The boundary expression must not be discarded.

Write $q=\partial_\nu u$, $\Pi=I-\nu\otimes\nu$, and
$a_\tau=\Pi A\nu$. Since the tangential derivatives of $u$ vanish,
$Du=q\nu$, $\Pi D^2u\,\nu=\nabla_\tau q$, and the tangential Hessian
is $q\mathcal S$, where $\mathcal S=D_\tau\nu$ is the shape operator
with this sign convention. Cancellation of the pure normal second
 derivative leaves precisely
\[
 (A:D^2u)q-q\nu\cdot D^2u A\nu
 =q^2\operatorname{tr}(\Pi A\Pi\mathcal S)
    +q\,a_\tau\cdot\nabla_\tau q.
\]
Tangential integration by parts on the closed boundary turns its integral into
\[
 \int_{\partial U}q^2
 \left(\operatorname{tr}(\Pi A\Pi\mathcal S)
              -\tfrac12\operatorname{div}_\tau a_\tau\right).
\]
The coefficient is bounded by the $C^1$ norm of $A$ and boundary curvature.
A tubular-neighborhood trace estimate gives, for any $\delta>0$,
\[
 \|Du\|_{L^2(\partial U)}^2
 \leq\delta\|D^2u\|_2^2+C_\delta\|Du\|_2^2.
\]
For completeness this follows by writing the boundary value of $|Du|^2$
as its value at an interior normal offset plus the normal integral of its
derivative, averaging offsets, and applying Young to $|Du||D^2u|$;
a finite collar cover has bounded Jacobians. Thus the boundary term costs
only the same small highest-order norm and a lower-order norm.
The drift and potential paired with $-\Delta u$ cost
$\delta\|D^2u\|_2^2+C_\delta\|u\|_{H^1}^2$.
It follows that
\[
 (Lu,-\Delta u)\geq c_1\|D^2u\|_2^2-C_1\|u\|_{H^1}^2.
\]
Finally $\|Du\|_2^2=-(u,\Delta u)$, so
$\|Du\|_2^2\leq\delta\|D^2u\|_2^2+C_\delta\|u\|_2^2$.
Absorb this term and add a sufficiently large multiple of $\|u\|_2^2$
to get \eqref{eq:7-h2-coercive}. Smooth zero-trace functions are dense
in $H^2\cap V$ on a smooth domain; all final volume pairings are continuous
in $H^2$, extending the estimate. Both source inequalities are the same
estimate. Our earlier reconstruction proofs avoid needing it inside the
Galerkin regularity passage, but it remains a useful independent result.
\end{proof}
\Needspace{55mm}
\PDEAnnotation{零边界不消去全部边界项}{$u=0$ 只消去切向一阶导数，法向导数 $q$ 可不为零。纯法向二阶项抵消后，仍有曲率项和 $q\nabla_\tau q$；须切向分部积分并用带小常数的迹估计吸收。}

\Needspace{30mm}
\paragraph{10. Telegraph equation uniqueness.}
For $u_{tt}+du_t-u_{xx}=f$ on $(0,1)$, with zero endpoint values and
two prescribed initial data, prove uniqueness for any real constant $d$.
\begin{proof}[Solution]
For a difference $w$ let $E=\frac12\int_0^1(w_t^2+w_x^2)$.
The endpoint term vanishes since $w_t=0$ there. Then
$E'=-d\int w_t^2\leq2|d|E$. Zero initial energy and Gronwall imply
$E=0$; zero endpoint values give $w=0$. For $d<0$ energy can grow,
but uniqueness still holds.
\end{proof}

\Needspace{30mm}
\paragraph{11. Clamped beam uniqueness.}
For $u_{tt}+u_{xxxx}=0$, $u=u_x=0$ at $0,1$, prove uniqueness.
\begin{proof}[Solution]
For a difference let $E=\frac12\int_0^1(w_t^2+w_{xx}^2)$.
Twice integrating gives
$\int w_{xxxx}w_t=\int w_{xx}w_{txx}$: the boundary terms are
$[w_{xxx}w_t-w_{xx}w_{tx}]_0^1=0$ since both clamped data are constant
in time. Thus $E'=0$, $E(0)=0$, and $w_{xx}=0$.
An affine function with zero clamped data is zero.
\end{proof}

\Needspace{30mm}
\paragraph{12. Resolvent identities.}
Establish the source identities (12)--(13) in \S7.4.1.
\begin{proof}[Solution]
On $D(G)$ the difference of the two operators is
$(\mu I-G)-(\lambda I-G)=(\mu-\lambda)I$.
Multiplication on the left by $R_\lambda$ and on the right by $R_\mu$
gives $R_\lambda-R_\mu=(\mu-\lambda)R_\lambda R_\mu$.
Interchange parameters to get commutation when they differ; equality
of parameters is automatic. The intermediate vectors are in $D(G)$,
so no multiplication of an undefined unbounded expression occurs.
\end{proof}

\Needspace{30mm}
\paragraph{13. Passing a closed generator through an integral.}
For $x\in D(G)$ justify
$G\int_0^\infty e^{-\lambda t}S(t)x\,dt
=\int_0^\infty e^{-\lambda t}GS(t)x\,dt$, $\lambda>0$.
\begin{proof}[Solution]
Equip $D(G)$ with $\|y\|_G=\|y\|+\|Gy\|$; closedness makes this a Banach
space. The integrand is graph-norm continuous because
$GS(t)x=S(t)Gx$, and its graph norm is at most
$e^{-\lambda t}(\|x\|+\|Gx\|)$.
Riemann sums on finite intervals therefore converge in graph norm, and
the tails converge there as well. Applying the continuous map
$G:(D(G),\|\cdot\|_G)\to X$ proves the identity.
\end{proof}
\Needspace{55mm}
\PDEAnnotation{图范数使积分合法}{只知道被积函数在 $X$ 中可积不足以交换无界算子。同时控制函数及其 $G$ 像，等于在完备图范数空间中积分。}

\Needspace{30mm}
\paragraph{14. Heat convolution on two different spaces.}
Let $S(t)g=\Phi_t*g$, $S(0)=I$. Show it is a contraction semigroup on
$L^2(\R^n)$, but not a strongly continuous semigroup on $L^\infty(\R^n)$.
\begin{proof}[Solution]
The Gaussian is nonnegative of mass one, so Young's inequality gives
$\|S(t)\|_{2\to2}\leq1$. Gaussian convolution gives
$S(t+s)=S(t)S(s)$. On $C_c^\infty$, approximate identities give
$S(t)g\to g$ in $L^2$; density and contraction extend this to all $L^2$.
On $L^\infty$ the same convolution still contracts the norm.
The failed axiom is strong continuity: take
$g=\mathbf1_{\{x_1>0\}}$. Its smoothed value tends to $1/2$ as
$x_1\downarrow0$, so for every $t>0$,
$\|S(t)g-g\|_\infty=1/2$ (in the essential-supremum sense).
The source's phrase ``not a contraction semigroup'' refers to this
continuity failure, not to failure of the norm inequality.
\end{proof}
\Needspace{55mm}
\PDEAnnotation{同一公式依赖所选空间}{热卷积在 $L^\infty$ 仍满足 $\|S(t)\|\leq1$，但阶跃函数阻止强连续。不能仅核对卷积律与范数，就宣布得到 $C_0$ 半群。}

\Needspace{30mm}
\paragraph{15. Invariance of powers of the generator domain.}
Define $D(G^k)=\{x\in D(G^{k-1}):G^{k-1}x\in D(G)\}$.
Show $S(t)D(G^k)\subset D(G^k)$.
\begin{proof}[Solution]
For $k=1$ this is Theorem~\ref{thm:7-generator}. Inductively, for
$x\in D(G^k)$, one has $S(t)x\in D(G^{k-1})$ and
$G^{k-1}S(t)x=S(t)G^{k-1}x\in D(G)$ by the $k=1$ case.
This proves invariance and also $G^kS(t)x=S(t)G^kx$.
\end{proof}

\Needspace{30mm}
\paragraph{16. Smooth spatial heat orbits from compactly supported data.}
For the Dirichlet heat semigroup and $g\in C_c^\infty(U)$, show
$S(t)g\in C^\infty(U)$ for every $t\geq0$.
\begin{proof}[Solution]
Here $G=\Delta$, $D(G)=H^2\cap V$.
Every $\Delta^jg$ is again compactly supported smooth, so
$g\in\bigcap_{k\geq1}D(G^k)$. Problem 15 preserves these domains.
Repeated elliptic regularity gives
$D(G^k)\subset H^{2k}(U)$: apply the elliptic estimate to
$Gu\in D(G^{k-1})$ inductively. For any spatial derivative order $r$,
choose $2k>r+n/2$ and use Sobolev embedding.
This proves the claim, including $t=0$. Abstract domain invariance alone
would not prove smoothing of arbitrary $L^2$ data; that stronger heat
property needs additional estimates.
\end{proof}
